\begin{trapbox}
	\begin{description}[style=nextline, leftmargin=2.8em, labelsep=0.5em, itemsep=0.35em]
		\item[\textbf{\textcolor{CTrap}{易错点一（焦点混淆）}}] 若椭圆焦点在 $y$ 轴，仍使用 $|PF|=a\pm e x_0$ 公式。
		\item[\textbf{\textcolor{CTrap}{正确理解（区分形式）}}] 焦点在 $y$ 轴时，焦半径公式应为 $|PF_1| = a + e y_0$（下焦点），$|PF_2| = a - e y_0$（上焦点）。
		\item[\textbf{\textcolor{CTrap}{错因分析（记混公式）}}] 记忆公式只记了 $x$ 轴的情形，忽略焦点位置对公式的影响。
	\end{description}

	\medskip
	\begin{description}[style=nextline, leftmargin=2.8em, labelsep=0.5em, itemsep=0.35em]
		\item[\textbf{\textcolor{CTrap}{易错点二（符号颠倒）}}] 将左、右焦半径公式写反：$|PF_1| = a - e x_0$，$|PF_2| = a + e x_0$。
		\item[\textbf{\textcolor{CTrap}{正确理解（左加右减）}}] 左焦半径（对应左焦点 $F_1(-c,0)$）为 $|PF_1|=a+e x_0$，右焦半径 $|PF_2|=a-e x_0$。
		\item[\textbf{\textcolor{CTrap}{错因分析（缺乏推导）}}] 只死记“左加右减”却不理解公式的由来，容易记反。
	\end{description}

	\medskip
	\begin{description}[style=nextline, leftmargin=2.8em, labelsep=0.5em, itemsep=0.35em]
		\item[\textbf{\textcolor{CTrap}{易错点三（忽略前提）}}] 对任意坐标 $x_0$ 直接使用焦半径公式，不验证点是否在椭圆上。
		\item[\textbf{\textcolor{CTrap}{正确理解（验证范围）}}] 公式成立的条件是点 $P(x_0,y_0)$ 在椭圆上，即 $x_0^2/a^2 + y_0^2/b^2 = 1$ 且 $x_0\in[-a,a]$。
		\item[\textbf{\textcolor{CTrap}{错因分析（遗忘条件）}}] 焦半径公式是椭圆上的点的性质，不能用于椭圆外的点。
	\end{description}


\medskip
\textbf{\textcolor{CTrap}{易错点四（第二定义忘了距离是非负的）}}
右准线 $x=\frac{a^2}{c}$ 对应的水平距离本质是
$\left|x_0-\frac{a^2}{c}\right|$，不能对任意点直接去绝对值。
只有因为\emph{点 $P$ 在椭圆上}，$|x_0|\le a<\frac{a^2}{c}$，
才能写成 $\frac{a^2}{c}-x_0$。

\medskip
\textbf{\textcolor{CTrap}{易错点五（从平方开方漏掉正负判断）}}
通常由 $|PF_2|^2=(a-ex_0)^2$ 先得到
$|PF_2|=|a-ex_0|$；再依据 $|x_0|\le a$、$e<1$，
证明 $a-ex_0\ge a-c>0$，才能删去绝对值。
同理 $a+ex_0\ge a-c>0$。

\medskip
\textbf{\textcolor{CTrap}{易错点六（用特殊位置代替证明）}}
在中心线 $x_0=0$ 和左右顶点 $x_0=\pm a$ 处验算都正确，
只能用来\emph{检查公式与判定符号}；不能仅凭这三处正确就断言整个椭圆上都成立。
正式证明仍需第一定义＋平方差、第二定义或距离公式。
\end{trapbox}
